\section{Updating the factorisation}
\label{sec:updates}

Each iteration changes the working set by one column, and the factors must follow.
Recomputing them from scratch costs a QR factorisation of $J\T A$, $O(nk^2)$; the
updates below cost $O(n^2)$ and $O(nk)$ respectively. Over the $O(n)$ iterations
of a run that is the difference between $O(n^3)$ and $O(n^4)$, and it is the whole
reason the factors are carried rather than rebuilt.

Both updates work the same way. The invariant~\eqref{eq:inv1} is preserved by
post-multiplying $J$ with any orthogonal $W$, since $(JW)\T G (JW) = W\T W = I$;
the update's task is to choose $W$ so that~\eqref{eq:inv2} is restored for the new
working set. Because $J' = JW$ gives ${J'}\T A = W\T(J\T A)$, the same $W\T$ acts on
the rows of $R$.

\subsection{Insertion}
\label{ssec:insert}

Let the working set of size $k$ be extended by the entering constraint, so
$A' = [\,A \ \ n_*\,]$. From~\eqref{eq:inv2} and $d = J\T n_*$,
\begin{equation}
  J\T A' = \begin{bmatrix} R & d_1 \\ 0 & d_2 \end{bmatrix},
  \label{eq:insert1}
\end{equation}
which fails to be triangular only in the $n-k$ entries of $d_2$ below its first.
Choose an orthogonal $W_2 \in \R^{(n-k)\times(n-k)}$ with
$W_2\T d_2 = \alpha e_1$ and set $W = \mathrm{diag}(I_k, W_2)$. Then
\begin{equation}
  (JW)\T A' = \begin{bmatrix} R & d_1 \\ 0 & W_2\T d_2\end{bmatrix}
  = \begin{bmatrix} R' \\ 0 \end{bmatrix},
  \qquad
  R' = \begin{bmatrix} R & d_1 \\ 0 & \alpha \end{bmatrix},
  \label{eq:insert2}
\end{equation}
upper triangular of order $k+1$, and the leading $k$ columns of $J$ are untouched
because $W$ acts as the identity on them.

\begin{proposition}[Full column rank is maintained]\label{prop:rank}
The insertion is performed only after a full step, which requires $t_2 < \infty$
and hence $z \neq 0$. Then $\alpha \neq 0$, so $R'$ is nonsingular and $A'$ has
full column rank $k+1$.
\end{proposition}

\begin{proof}
$|\alpha| = \norm{W_2\T d_2} = \norm{d_2}$, which is nonzero iff $z \neq 0$ by
Lemma~\ref{lem:directions}(ii). $R'$ is triangular with diagonal that of $R$
extended by $\alpha$, so it is nonsingular, and ${R'}\T R' = {A'}\T G^{-1}A'$ is then
positive definite, which forces $A'$ to have full column rank.
\end{proof}

The proposition is worth pausing on: the algorithm never has to test for
linear dependence among the working-set normals. The step rules exclude it. A
constraint whose normal lies in the span of those already held has $z = 0$ by
Lemma~\ref{lem:directions}(iv), so its step is limited by $t_1$ and it can only
be added after enough partial steps have removed the dependence. This is the
structural advantage that Section~\ref{sec:pdas} gives up.

\paragraph{The choice of $W_2$.} Any orthogonal $W_2$ with
$W_2\T d_2 \in \R e_1$ will do, and by Proposition~\ref{prop:invariance} the
choice does not affect the iterates. The reference implementation uses a chain of
$n-k-1$ Givens rotations, one per trailing component. A single Householder
reflection achieves the same reduction:
\begin{equation}
  W_2 = I - \frac{2}{\beta}\, v v\T,
  \qquad v = d_2 - \alpha e_1,
  \qquad \beta = v\T v,
  \qquad \alpha = \pm\norm{d_2}.
  \label{eq:householder}
\end{equation}
That $W_2 d_2 = \alpha e_1$ is the standard Householder identity, using
$\beta = 2\,v\T d_2$. The reflection is symmetric, so it applies to the columns of
$J$ and to $d_2$ as one operation, and applying it to a block is a
matrix--vector product followed by a rank-one update --- two operations of
$O(n(n-k))$, where the Givens chain is $n-k-1$ operations of $O(n)$. The
arithmetic is comparable; the number of operations is not.

\begin{remark}[The sign of $\alpha$, and cancellation]\label{rem:sign}
Numerical analysis prescribes $\alpha = -\mathrm{sign}((d_2)_1)\norm{d_2}$, which
makes $v_1 = (d_2)_1 - \alpha$ a sum of like-signed terms. The implementation
takes the opposite sign, $\alpha = \mathrm{sign}((d_2)_1)\norm{d_2}$, so as to
reproduce the sign convention of the reference's Givens chain --- and thereby walks
into the cancellation the prescription exists to avoid, severely when $d_2$ is
already close to a positive multiple of $e_1$. The cure is algebraic rather than a
change of convention: with $\tau = \norm{(d_2)_{2:}}^2$ and
$\nu = \norm{d_2} = \sqrt{(d_2)_1^2 + \tau}$,
\begin{equation}
  v_1 = (d_2)_1 - \mathrm{sign}((d_2)_1)\,\nu
      = \frac{(d_2)_1^2 - \nu^2}{(d_2)_1 + \mathrm{sign}((d_2)_1)\nu}
      = \frac{-\mathrm{sign}((d_2)_1)\,\tau}{|(d_2)_1| + \nu},
  \label{eq:nocancel}
\end{equation}
in which every operation combines like-signed quantities. The remaining entries
of $v$ are those of $d_2$, and $\beta = \tau + v_1^2$. By
Proposition~\ref{prop:invariance} the sign convention is anyway immaterial to the
iterates; \eqref{eq:nocancel} is what makes it immaterial to the accuracy as
well. See~\cite{golub2013,higham2002} for the standard treatment.
\end{remark}

\begin{remark}[Why the reflections cannot be blocked]\label{rem:wy}
A sequence of Householder reflections is normally applied as a block. The WY
representation of Bischof and Van Loan~\cite{bischof1987}, and its storage-efficient
form due to Schreiber and Van Loan~\cite{schreiber1989}, write a product
$W_1 W_2 \cdots W_j$ as $I + WY\T$, so that $j$ reflections reach a matrix in one
level-3 update rather than $j$ level-2 ones. That is what makes blocked QR
factorisation fast, and it is the first thing to reach for here.

It is not available, for a structural reason rather than an implementation one.
Blocking requires the $j$ reflections to be known before any is applied. In this
iteration the reflection that inserts a constraint is determined by
$d = J\T n_*$, which depends on $J$ \emph{after} the previous insertion; and
between any two insertions the solver reads $z$ and $r$ off the updated factors to
decide which constraint enters next, and whether it enters at all. The updates are
interleaved with reads of what they produce, so there is no window in which two
reflections are simultaneously known and unapplied.

This is the same obstruction as in the deletion chase below, and it is why the
method stays a level-2 computation no matter how it is coded, which is worth
knowing before reaching for a blocked update. A variant
that inserted several constraints per iteration would open the window; the dual
method's step rules, which admit one constraint at a time by construction, close
it.
\end{remark}

\subsection{Deletion}
\label{ssec:delete}

Let the constraint in position $p$ of the working set be removed. Deleting column
$p$ of $A$ deletes column $p$ of $R$ and shifts columns $p+1,\dots,k$ one place
left, each retaining its own entries. The result is upper triangular except for
one subdiagonal: writing $\tilde R$ for the shifted $k \times (k-1)$ matrix,
$\tilde R_{ij} = 0$ for $i > j+1$, with the offending entries
$\tilde R_{j+1,j}$ for $j = p,\dots,k-1$.

These entries lie in distinct columns, and no single reflection annihilates
components of distinct vectors. What restores the shape is a \emph{chase} of
$k-p$ plane transformations: for $j = p,\dots,k-1$ in order, a $2\times2$
orthogonal acting on rows $j$ and $j+1$ annihilates $\tilde R_{j+1,j}$, and the
same transformation is applied to columns $j$ and $j+1$ of $J$. The chase is
inherently sequential --- the parameters of each transformation are read from
entries the previous one wrote --- which is why no batched form of it exists.

The implementation takes the $2 \times 2$ block to be the symmetric reflection
\begin{equation}
  W_j^{(2)} = \begin{bmatrix} c & s \\ s & -c \end{bmatrix},
  \qquad c = \frac{x}{h},\quad s = \frac{y}{h},
  \qquad h = \mathrm{sign}(x)\sqrt{x^2+y^2},
  \label{eq:refl2}
\end{equation}
where $x = \tilde R_{jj}$ and $y = \tilde R_{j+1,j}$, so that
$W_j^{(2)}(x,y)\T = (h,0)\T$. Symmetry is the point. The update rule of this
section applies $W$ to the columns of $J$ and $W\T$ to the rows of $R$; a
symmetric block makes those the same matrix, so one $2\times2$ serves both sides.
A Givens rotation, which is what BLAS offers, is not symmetric: it agrees
with~\eqref{eq:refl2} on the first output and negates the second, so it requires
its transpose on the $R$ side. By Proposition~\ref{prop:invariance} either choice
gives the same iterates, the difference being a sign matrix.

\begin{remark}[The degenerate step]
When $x = 0$ the formula~\eqref{eq:refl2} gives $c = 0$, $s = \pm1$, and the
transformation reduces to an exchange of the two rows up to sign. The
implementation performs a plain exchange, $\left[\begin{smallmatrix} 0 & 1 \\ 1 &
0\end{smallmatrix}\right]$, which is orthogonal and symmetric and annihilates $y$
just as well; it differs from~\eqref{eq:refl2} by a sign when $y < 0$, and
Proposition~\ref{prop:invariance} says that is free.
\end{remark}

\begin{remark}[Where the storage cost lands]
$R$ is held as packed columns --- entry $(i,j)$, $i \le j$, at offset
$j(j+1)/2 + i$ --- so its leading $k \times k$ block, which the solve for $r$ reads
once per iteration, is a contiguous run at every $k$. Insertion writes one column,
contiguous in that layout; deletion mixes two \emph{rows} across a range of
columns, which is a strided gather because column offsets grow with the index. The
asymmetry is deliberate, and it is the right way round. It is a property of the
layout rather than of the derivation, to which the mathematics is indifferent.
\end{remark}
